
\documentclass{article}
\usepackage{colortbl}
\usepackage{makecell}
\usepackage{multirow}
\usepackage{supertabular}

\begin{document}

\newcounter{utterance}

\centering \large Interaction Transcript for game `ifeval', experiment `ifeval', episode 367 with Qwen3.5{-}27B.
\vspace{24pt}

{ \footnotesize  \setcounter{utterance}{1}
\setlength{\tabcolsep}{0pt}
\begin{supertabular}{c@{$\;$}|p{.15\linewidth}@{}p{.15\linewidth}p{.15\linewidth}p{.15\linewidth}p{.15\linewidth}p{.15\linewidth}}
    \# & \multicolumn{2}{c}{Player} && \multicolumn{2}{c}{Game Master} \\
    \hline

    \theutterance \stepcounter{utterance}  
    & & & \multicolumn{4}{p{0.6\linewidth}}{
        \cellcolor[rgb]{0.9,0.9,0.9}{
            \makecell[{{p{\linewidth}}}]{
                \texttt{\tiny{[P1$\langle$GM]}}
                \texttt{Expand the riddle into a story with a funny tone:} \\
\\ 
\texttt{What can you catch but not throw?} \\
\texttt{A cold} \\
\\ 
\texttt{Use * to highlight at least 2 sections in your text. For example: *this is a highlighted text section*.} \\
            }
        }
    }
     \\ \\

    \theutterance \stepcounter{utterance}  
    & \multicolumn{4}{p{0.6\linewidth}}{
        \cellcolor[rgb]{0.561,0.737,0.561}{
            \makecell[{{p{\linewidth}}}]{
                \texttt{\tiny{[P1$\rangle$GM]}}
                \texttt{Barnaby "The Cannon" McGee was a man who took the concept of "throwing" very seriously. Whether it was a football, a tantrum, or a half{-}eaten sandwich at a stray cat, Barnaby believed that if he could grab it, he could launch it with the force of a small rocket. One particularly chilly Tuesday, however, Barnaby encountered an opponent that would test the very limits of his physics{-}defying philosophy.} \\
\\ 
\texttt{It started with a sneeze that sounded like a deflating bagpipe. Then came the sniffles, the watery eyes, and the general feeling that his nose had been replaced by a wet sponge. Barnaby looked at his reflection in the mirror, sniffled loudly, and declared, "I've been caught! But don't worry, folks, because I'm about to *throw this thing right back at the universe*!"} \\
\\ 
\texttt{He stood in the center of his living room, wound up his arm like a professional pitcher, and hurled an invisible ball of germs toward the ceiling fan. "Take that, you pesky virus!" he shouted. The fan just whirred on, indifferent to his athletic display. He tried again, aiming for the cat. "Catch this, Whiskers!" he yelled, making a dramatic throwing motion. The cat merely blinked, licked its paw, and walked away, clearly unimpressed by Barnaby's attempt to toss a non{-}existent object.} \\
\\ 
\texttt{Barnaby was confused. He had caught the cold, sure, but the laws of physics seemed to have taken a coffee break. He tried to throw it at the door, at the toaster, and even at his own reflection, but every time he made the motion, the cold just stayed right where it was: stuck in his sinuses, mocking his lack of coordination.} \\
\\ 
\texttt{*It turns out that while you can catch a cold with the ease of a fly catching a fly, you absolutely cannot throw it like a baseball.* The more Barnaby tried to launch his illness into the stratosphere, the more he realized that some things are meant to be endured, not lobbed over a fence.} \\
\\ 
\texttt{By the time he finally accepted his fate, he was wrapped in three blankets, drinking tea that tasted like warm dishwater, and sneezing into a pillow. He looked at his empty hands and sighed. "Well," he muttered, "at least I didn't throw it at the cat."} \\
\\ 
\texttt{*The moral of the story is simple: never try to play catch with a virus, because the only thing you'll end up throwing is your dignity.* Barnaby learned that day that some things are strictly for catching, and if you try to throw them, you'll just end up looking silly while your nose runs down your chin.} \\
            }
        }
    }
    & & \\ \\

    \theutterance \stepcounter{utterance}  
    & & & \multicolumn{2}{p{0.3\linewidth}}{
        \cellcolor[rgb]{0.9,0.9,0.9}{
            \makecell[{{p{\linewidth}}}]{
                \texttt{\tiny{[GM$|$GM]}}
                {'detectable_format:number_highlighted_sections': {'num_highlights': 2}}
            }
        }
    }
    & & \\ \\

    \theutterance \stepcounter{utterance}  
    & & & \multicolumn{2}{p{0.3\linewidth}}{
        \cellcolor[rgb]{0.9,0.9,0.9}{
            \makecell[{{p{\linewidth}}}]{
                \texttt{\tiny{[GM$|$GM]}}
                \texttt{game\_result = WIN} \\
            }
        }
    }
    & & \\ \\

\end{supertabular}
}

\end{document}
